\documentclass[11pt,a4paper]{article}
\usepackage[T1]{fontenc}
\usepackage[utf8]{inputenc}
\usepackage{lmodern,amsmath,amssymb,amsthm,mathtools,microtype,mathrsfs}
\usepackage[margin=25mm,top=23mm,bottom=25mm]{geometry}
\usepackage{booktabs,array,longtable,enumitem,xcolor,fancyhdr}
\usepackage[colorlinks=true,linkcolor=blue!45!black,urlcolor=blue!45!black,citecolor=blue!45!black]{hyperref}
\hypersetup{pdftitle={Reading the Alpoge--Claude RH and Dirichlet-GRH proof: working version 0.37},pdfauthor={ChatGPT},pdfsubject={Partial human-readable reconstruction with explicit arithmetic inputs},pdfkeywords={Riemann hypothesis, Lean, theta lifts, prime sums, Arakelov slopes}}
\setlength{\parindent}{0pt}
\setlength{\parskip}{5pt plus 1pt}
\setlist{topsep=3pt,itemsep=3pt,leftmargin=*}
\setlength{\headheight}{14pt}
\pagestyle{fancy}\fancyhf{}
\fancyhead[L]{\small RH and Dirichlet GRH: a reconstruction}
\fancyhead[R]{\small Working version 0.37}
\fancyfoot[C]{\thepage}
\renewcommand{\headrulewidth}{0.3pt}
\numberwithin{equation}{section}
\newtheorem{theorem}{Theorem}[section]
\newtheorem{proposition}[theorem]{Proposition}
\newtheorem{lemma}[theorem]{Lemma}
\newtheorem{corollary}[theorem]{Corollary}
\theoremstyle{remark}\newtheorem{remark}[theorem]{Remark}
\newcommand{\Rea}{\operatorname{Re}}
\newcommand{\HS}{\mathrm{HS}}
\newcommand{\tr}{\operatorname{tr}}
\newcommand{\codim}{\operatorname{codim}}
\newcommand{\covol}{\operatorname{covol}}
\newcommand{\OK}{\mathcal O_K}
\newcommand{\src}[2]{\href{https://github.com/CocoSgt/riemann-hypothesis-lean/blob/7cbf7a0c32fca1badf05931879a37c3f8a2b03b4/src/#1}{\texttt{#2}}}
\newcommand{\status}[1]{\textbf{Status:} #1}
\begin{document}
\section{Bounding the initial theta function before its Hilbert realization}
\label{sec:seed-boundedness}

This chapter closes the initial boundedness input left in SPR25.  It
works on the actual function supplied by the original theta integral.
The finite-volume and finite-class inputs are FV1--FV11.  Central
finiteness is proved directly from the original extremal equations,
before any target Hilbert vector is constructed.  The only analytic
import is a precisely matched rapid-decay theorem.  The universal
source predicate \texttt{Published.\allowbreak Cusp\allowbreak Rapid\allowbreak Decay}, with its stronger
global quotient-height conclusion, is not required for this route.

Throughout, put \(F=D.\mathrm{host}.F\),
\(E=D.\mathrm{host}.E\), \(\sigma=D.\mathrm{sp.cm.sig}\),
and \(J=V.\mathrm{gram}\).  Use the same original
\(G=\mathrm{GA}(D,V)\), rational subgroup \(\Gamma=\mathrm{GQ}(V)\),
infinite family \(\theta\), and real and finite factors as in HCU14,
RQC1 and FV4.  Thus \(G=G_fG_\infty=G_f\times G_\infty\), with
commuting factors, and
\[
 G_\infty\simeq\prod_{\tau\in\mathrm{Tau}}U(J_\tau),\qquad
 \mathbf G=\operatorname{Res}_{F/\mathbb Q}U(J),\qquad
 G_\infty=\mathbf G(\mathbb R).
 \tag{SBD.1}
\]
The family \(\theta\) is assumed throughout to satisfy the original
\(\mathrm{IsInfFam}(\theta)\) predicate.
Here and below an identification means the explicit matrix and
topological identification of FV4, not a replacement of the source
group.  The same conventions govern the derivative \(R_X\): its
curve at \(g\) is \(g\exp(tX)\).  Operator products act on a function
from right to left.  No target quotient measure is used in SBD1--SBD13.

% BEGIN SBD1
\paragraph{SBD1: the exact smooth function module and its Lie action.}
\label{claim:SBD1}
Let \(\mathcal C\) be the complex vector space of functions on \(G\)
whose restriction to every coset \(xG_\infty\) is smooth.  For
\(X\in\mathfrak u(J_\tau)\), the source operator
\(R_X=\mathrm{Rder}(\theta,V,\tau,X)\) preserves \(\mathcal C\).
Its complex extension is literally \(\mathrm{RderC}\).  It gives a
representation of
\[
 \mathfrak g_{\mathbb C}
   =\bigoplus_\tau\mathfrak{gl}_{n}(\mathbb C),
 \qquad n=\operatorname{rank}V,
\]
and hence of its universal enveloping algebra on \(\mathcal C\).
The source curve-product smoothness implies membership in
\(\mathcal C\).  These assertions require no boundedness, Hilbert
space, quotient norm or formal adjoint.

\emph{Proof.}
ASB1--ASB3 identify the original real matrix groups and their
one-place exponential curves.  Choose a real basis of their product
Lie algebra.  A product of the corresponding exponentials is a
smooth coordinate chart around the identity; its differential there
is the identity in these coordinates.  The original smoothness
condition applies to this particular list and every left translate.
It therefore makes each restriction a smooth function in charts.
Conversely any smooth restriction has the original curve-product
smoothness, because finite multiplication and exponential are smooth.

On a real matrix group the vector field defining \(R_X\) has value
\(gX\) at \(g\).  For a smooth ambient extension in a local chart,
the two second derivatives in the commutator cancel by symmetry,
leaving \(g(XY-YX)\).  Thus
\[
 [R_X,R_Y]=R_{[X,Y]}.
 \tag{SBD.2}
\]
This calculation is intrinsic, since each vector field is tangent
to the group.  Directional differentiation is real linear in the
direction, complex linear in the function, and again smooth.
Different product factors commute, so their operators commute too.

For the actual nonsingular Hermitian matrix \(J_\tau\), set
\(c(Z)=-J_\tau^{-1}Z^*J_\tau\).  This is an antilinear involution
with fixed space \(\mathfrak u(J_\tau)\), and
\[
 Z=P+iQ,\qquad P=\tfrac12(Z+c(Z)),\quad
 Q=\tfrac1{2i}(Z-c(Z)).
 \tag{SBD.3}
\]
Consequently complexification is \(R_P+iR_Q\), exactly the two
summands in the definition of \texttt{Rder\allowbreak C}.  Real linearity and
(SBD.2) imply the complex Lie relations.  The tensor algebra map
defined by ordered operator products kills
\(X\otimes Y-Y\otimes X-[X,Y]\); factoring this map gives the
enveloping-algebra representation.  In particular
\[
 [R_{E_{ab}},R_{E_{cd}}]
   =\delta_{bc}R_{E_{ad}}-\delta_{da}R_{E_{cb}}.
 \tag{SBD.4}
\]
There is no sign reversal or change in word order.
% END SBD1

% BEGIN SBD2
\paragraph{SBD2: highest-weight equations for the initial extremal function.}
\label{claim:SBD2}
Retain the original \(M,\theta,w\), satisfying
\(\mathrm{HolShape}\), \(\mathrm{IsInfFam}\) and \(\mathrm{WtIs}\).
Let \(U\) be the original lowest-type multilinear map with smooth
values, compact equivariance \(\mathrm{KEquiv}\), and the original
lowering equations.  If \(e\) is an original extremal tuple, put
\(\phi=U(e)\) and \(x_\tau=\mathrm{x0Mat}(M,\tau)\).  Then
\[
 \begin{aligned}
 R_{x_\tau E_{ab}x_\tau^{-1}}\phi&=0 &&(a<b),\\
 R_{x_\tau E_{aa}x_\tau^{-1}}\phi&=w_{\tau a}\phi.
 \end{aligned}
 \tag{SBD.5}
\]
These are identities of the initial functions on all of \(G\).

\emph{Proof.}
Use the proof of HCU19, whose intermediate conclusion is the
literal \(\mathrm{KilledFn}\) identity for
\(u\mapsto U(e)(gL(u))\), for every \(g\).  The compact lowest
types are finite dimensional by HCU13.  Differentiating compact
equivariance is therefore ordinary finite-dimensional differentiation
of evaluation on those types.  The compact upper units kill the
extremal vector and the diagonal units give the original weight.
The other upper units, in the \(x_\tau\) ordering, are exactly the
original lowering directions and kill every value of \(U\).
HCU14 identifies their curves with the source one-place curves.
Evaluating this \(\mathrm{KilledFn}\) identity at the identity real
tuple, for each \(g\), gives (SBD.5).  This uses the intermediate
equations in HCU19, not an unproved converse to \(\mathrm{HolAd}\).
No target Hilbert realization occurs.  The zero function is allowed.
% END SBD2

% BEGIN SBD3
\paragraph{SBD3: every central enveloping operator acts by a scalar.}
\label{claim:SBD3}
Let a module for a finite direct sum of complex matrix Lie algebras
contain a vector \(v\) satisfying the upper-triangular equations
(SBD.5), in specified invertible frames.  For every element \(z\)
of the center of the full universal enveloping algebra, there is a
scalar \(c_z\), determined by \(z\), the frames and the weights,
such that \(zv=c_zv\).  The statement does not require
irreducibility, a Hilbert norm, an injective module realization, or
linear independence of ordered monomials.

\emph{Proof.}
Conjugating the matrix units by the specified frames preserves
(SBD.4); include this conjugation in their names.  Order all the
strictly lower units first, all diagonal units next, and all
strictly upper units last, choosing fixed orders within these lists.
Every word is a linear combination of ordered words.  Indeed swap
an adjacent inversion by \(AB=BA+[A,B]\).  The swapped term has
fewer inversions at the same length, whereas the bracket terms in
(SBD.4) have smaller length.  Induction first on length and then on
the inversion count terminates.

Give \(E^{(\tau)}_{ab}\) weight
\(\epsilon_{\tau a}-\epsilon_{\tau b}\).  Every word is an
eigenvector for the adjoint action of the diagonal Cartan algebra,
with the sum of its letter weights.  We may express \(z\) using
only weight-zero words, without assuming a monomial basis.  To see
this, take any finite expression for \(z\) as words.  Choose a
diagonal \(H\) on which none of its finitely many nonzero word
weights vanishes.  Such an \(H\) exists because finitely many proper
complex hyperplanes cannot cover the Cartan space.  If \(S\) is
the finite set of its nonzero eigenvalues, apply
\[
 p(\operatorname{ad}H),\qquad
 p(T)=\prod_{a\in S}(1-T/a).
 \tag{SBD.6}
\]
It fixes \(z\), since \([H,z]=0\), and removes every nonzero-weight
word.  The adjacent reordering above preserves total weight in
each surviving bracket term.  Thus \(z\) is a linear combination
of ordered weight-zero words.

An ordered word containing a strictly upper unit kills \(v\),
because its rightmost upper unit is applied first.  If it contains
no upper unit but does contain a lower unit, its weight cannot be
zero: the linear functional sending
\(\epsilon_{\tau a}-\epsilon_{\tau b}\) to \(b-a\), summed over
the factors, is strictly negative on lower units and zero on
diagonal ones.  Therefore every surviving word is a product of
diagonal units.  It acts by the corresponding product of the
weights.  Summing gives \(c_z\).  This also proves that the scalar
may be chosen independently of the particular vector or module
with these equations.  If \(v\ne0\), uniqueness shows that
\(z\mapsto c_z\) is a character of the center.  If \(v=0\), the
conclusion is immediate.  Rank zero gives the scalar algebra and
is immediate as well.
% END SBD3

% BEGIN SBD4
\paragraph{SBD4: the literal cyclic operators and full center finiteness.}
\label{claim:SBD4}
For the initial function in SBD2, the one-dimensional space
\(W=\mathbb C\phi\) is stable both under the full center of
\(U(\mathfrak g_{\mathbb C})\) and under every original
\(\mathrm{gelf}(\theta,V,\tau,k)\).  Hence it supplies the literal
\texttt{Is\allowbreak Aut\allowbreak Fun\allowbreak Sm.\allowbreak zfin} field and the full center-finiteness
hypothesis used below.  For a fixed original model, the cyclic
scalars depend only on its weights, not on the theta parameter or
source vector.

\emph{Proof.}
Full center stability follows from SBD1--SBD3.  For completeness,
let \(\mathcal E=(E_{ab})\) have entries in the enveloping algebra
of one factor, and write \(Z_m=\operatorname{tr}(\mathcal E^m)\)
for \(m\ge1\).  The matrix-unit relations imply, by induction on
\(m\),
\[
 [E_{ab},(\mathcal E^m)_{ij}]
   =\delta_{bi}(\mathcal E^m)_{aj}
      -\delta_{aj}(\mathcal E^m)_{ib}.
 \tag{SBD.7}
\]
In the induction step, apply the derivation rule to
\(\sum_l(\mathcal E^m)_{il}E_{lj}\); the two middle sums cancel
with the same operator order.  Summing (SBD.7) over \(i=j\)
gives \([E_{ab},Z_m]=0\).  Thus \(Z_m\) is central in that
factor and in the full direct-sum enveloping algebra.

Expanding its trace gives the ordered closed paths
\[
 Z_{k+1}=\sum_{i_0,\ldots,i_k}
 E_{i_0i_1}E_{i_1i_2}\cdots E_{i_ki_0}.
 \tag{SBD.8}
\]
Under SBD1 this is exactly the source right-fold operator
\(\mathrm{gelf}(\theta,V,\tau,k)\): its length is \(k+1\), with
no factorial and no additional sign.  Alternatively, conjugating
the generators by \(x_\tau\) leaves this trace unchanged by the
scalar-matrix trace calculation in RQC17.  Reordering its
weight-zero words as in SBD3 shows that its scalar is an
integer-coefficient polynomial in the original weights.  For a
zero seed, \(W=0\) and all requirements remain true.  This proof
does not import the post-Hilbert scalar theorem SPR18 or a theorem
that the cyclic elements generate the center.
% END SBD4

% BEGIN SBD5
\paragraph{SBD5: the original growth becomes the required real moderate growth.}
\label{claim:SBD5}
Let \(\phi\) be an original automorphic function with source
smoothness and \(\mathrm{IsUnifModGrowth}\).  Fix a compact open
finite-level subgroup \(K\) fixing \(\phi\) on the right, and
\(s\in G_f\).  Put
\[
 F_s(a)=\phi(sa),\qquad
 \Gamma_s=\{\gamma_\infty:\gamma\in U(J)(F),\
                         \gamma_f\in sKs^{-1}\}.
 \tag{SBD.9}
\]
Then \(F_s\) is a smooth \(\Gamma_s\)-invariant function on the
actual \(G_\infty\), where \(\Gamma_s\) is arithmetic.  Its
source derivative-growth estimates imply moderate growth in a
faithful linear norm; the same source exponent works for every
ordered real derivative, with constants allowed to depend on
the word and \(s\).

\emph{Proof.}
FV10 proves the arithmetic assertion using the actual rational
matrices and the compact open \(sKs^{-1}\).  If
\(\gamma_f\in sKs^{-1}\), then
\[
 (s,\gamma_\infty a)=\gamma(s,a)
                  (s^{-1}\gamma_f^{-1}s,1).
\]
Rational left invariance and right \(K\)-invariance give
\(F_s(\gamma_\infty a)=F_s(a)\).  SBD1 gives smoothness.

Realify the faithful rational representation of
\(\operatorname{Res}_{F/\mathbb Q}U(J)\) on \(E^n\), and take
its direct sum with its dual and a trivial one-dimensional
representation.  Let \(\|a\|_{\mathrm{lin}}\) be the operator
norm.  It is submultiplicative and at least one.  In the actual
CM-type coordinates the original representation is the direct
sum of the complex matrices \(a_\tau\), regarded as real matrices;
the dual controls their inverse matrices.  Equivalence under the
one fixed change of real basis therefore gives, for \(n>0\),
\[
 \mathrm{normInf}(\theta,V,a)
   =\sum_{\tau,i,j}(|(a_\tau)_{ij}|+|(a_\tau^{-1})_{ij}|)
   \le C_0\|a\|_{\mathrm{lin}}.
 \tag{SBD.10}
\]
For example, in an orthonormal block basis each entry is bounded
by its block operator norm; a fixed basis change contributes only
a constant.  The same \(C_0\) works for every \(a\).
The source growth estimate at \(x=s\), and the empty word, is
the required moderate-growth inequality.  For a nonempty word,
SBD1 identifies the operators exactly and the same comparison
applies.  A finite linear combination of words retains this
exponent after increasing its constant.  Rank zero is the
one-point group and all bounds are immediate.

An original open subgroup in \texttt{Is\allowbreak Aut\allowbreak Fun} contains such a
finite compact open \(K\): intersect its finite slice with an
integral compact open from FV4.  The intersection is a closed
subgroup of that compact group and is open.  Thus no stronger
level hypothesis has been introduced.
% END SBD5

% BEGIN SBD6
\paragraph{SBD6: a finite-level fundamental domain for the additive adeles.}
\label{claim:SBD6}
Let \(B\subset\mathbb A_{F,f}\) be a compact open additive
subgroup and put \(\Lambda_B=F\cap B\), with its real embedding
in \(F_\infty\).  Then
\(\mathbb A_{F,f}=F+B\), \(\Lambda_B\) is a full real lattice,
and, for a half-open lattice parallelepiped \(Q_B\),
\[
             Q_B\times B
 \quad\hbox{is a measurable fundamental domain for }
             F\backslash\mathbb A_F.
 \tag{SBD.11}
\]
Its closure is compact and its product Haar measure is finite
and positive.

\emph{Proof.}
There are positive integers \(b,c\) with
\(b\widehat{\mathcal O}_F\subset B\subset
c^{-1}\widehat{\mathcal O}_F\).  The first inclusion is the
neighborhood basis in the finite restricted product; the second
follows by compactness, since the sets
\((m!)^{-1}\widehat{\mathcal O}_F\) form an increasing open
cover.  For \(y\in\mathbb A_{F,f}\), clear its finitely many
denominators by an integer \(d>0\).  The Chinese remainder
theorem chooses \(a\in\mathcal O_F\) whose residues agree with
those of \(dy\) modulo \(db\widehat{\mathcal O}_F\).
Then \(y-a/d\in b\widehat{\mathcal O}_F\subset B\).
This proves the claimed additive approximation.

The elementary local-to-global integrality test of ASB7 gives
\[
 b\mathcal O_F\subset\Lambda_B\subset c^{-1}\mathcal O_F.
\]
An integral basis and the nonsingular embedding matrix, as in
ASB9, make \(\mathcal O_F\) a full lattice in \(F_\infty\).
The quotient \(c^{-1}\mathcal O_F/b\mathcal O_F\) is finite,
so \(\Lambda_B\) is a full lattice as well.

Subtract a rational element to move the finite part of any
adele into \(B\); then subtract an element of \(\Lambda_B\)
to move its real part into \(Q_B\).  This proves coverage.
Two points of \(Q_B\times B\) differing by a diagonal rational
element differ by an element of \(\Lambda_B\), so the half-open
parallelepiped forces equality.  Alternatively a closed
parallelepiped gives uniqueness outside its measure-zero faces.
The product is Borel and has the stated measure properties,
since \(B\) is nonempty compact open and \(Q_B\) is a full
lattice parallelepiped.
% END SBD6

% BEGIN SBD7
\paragraph{SBD7: the original transvection subgroup is the actual additive line.}
\label{claim:SBD7}
Let \(0\ne\xi\in E^n\) be an original isotropic vector,
\(\xi^*J\xi=0\).  Choose \(\eta\in E^\times\) with
\(\sigma(\eta)=-\eta\), and put
\[
 T_\xi=\xi\xi^*J,\qquad
 z_\xi(a)=1+\eta a T_\xi.
 \tag{SBD.12}
\]
Then \(a\mapsto z_\xi(a)\) identifies the additive groups
\(F\), \(\mathbb A_{F,f}\), \(F_\infty\), and \(\mathbb A_F\)
with the rational, finite, real, and full points of a closed
unipotent subgroup.  On the full group it is exactly the source
\(\mathrm{ZI}(V,\xi)\), and
\[
       \mathrm{ZI}(V,\xi)\cap\Gamma=z_\xi(F).
 \tag{SBD.13}
\]

\emph{Proof.}
Take \(\eta=e-\sigma(e)\ne0\); the fixed field is \(F\), so
\(E=F\oplus F\eta\).  Since \(T_\xi^2=0\),
\(z_\xi(a)z_\xi(b)=z_\xi(a+b)\) and the inverse is
\(z_\xi(-a)\).  Also \(T_\xi^*J=JT_\xi\).  Hence
\[
 (1+uT_\xi)^*J(1+uT_\xi)
      =J+(u+\sigma(u))JT_\xi.
 \tag{SBD.14}
\]
The nonzero matrix \(JT_\xi\) has a nonzero entry in \(E\),
which is a unit after the original adelic scalar inclusion.
Thus this matrix is unitary exactly when
\(u+\sigma(u)=0\), equivalently \(u=\eta a\) for a unique
\(a\) in the corresponding scalar algebra over \(F\).

The generating set in the source definition of \(\mathrm{ZI}\)
is therefore already this subgroup.  Its group-generated closure
adds no elements.  A nonzero entry of \(\eta T_\xi\) recovers
\(a\) continuously from \(z_\xi(a)-1\); the other entries are
fixed linear multiples of it.  This proves the topological
identification and closedness.  If the resulting matrix is
rational, this same entry shows that \(u\in E\), and the
anti-invariance equation gives \(a\in F\).  Conversely rational
parameters give original rational unitary matrices.  This proves
(SBD.13), including its subgroup and topology conventions.
% END SBD7

% BEGIN SBD8
\paragraph{SBD8: original adelic zero averages give the exact real zero averages.}
\label{claim:SBD8}
Suppose \(\phi\) is continuous, rationally left invariant,
right \(K\)-invariant, and satisfies the original
\(\mathrm{ZCusp}(V,\phi)\).  For every \(s\in G_f\), every
original isotropic \(\xi\), and every \(a\in G_\infty\),
\[
 \int_{(\Gamma_s\cap Z_\xi(\mathbb R))\backslash
                  Z_\xi(\mathbb R)} F_s(za)\,dz=0.
 \tag{SBD.15}
\]
Any Haar normalization of this compact real quotient is allowed.

\emph{Proof.}
In the additive parameter of SBD7 set
\[
 B=\{b\in\mathbb A_{F,f}:z_\xi(b)\in sKs^{-1}\}.
\]
This is compact open: openness follows from continuity, and
compactness from the closed embedding in SBD7 and compactness of
\(sKs^{-1}\).  Moreover \(\Gamma_s\cap Z_\xi(\mathbb R)\)
has exactly the parameter lattice \(\Lambda_B\).
Use the product domain \(Q_B\times B\) of SBD6, transported to
the \emph{same} original \(\mathrm{ZI}\), and transport product
Haar measure there.  The universal measure/domain quantifiers in
\(\mathrm{ZCusp}\) allow this choice.  At the argument \(sa\),
the original zero integral becomes
\[
 0=\int_{Q_B}\int_B
       \phi\bigl(z_\xi(u)_\infty z_\xi(b)_fsa\bigr)\,db\,du
   =\operatorname{vol}(B)
       \int_{Q_B}F_s\bigl(z_\xi(u)_\infty a\bigr)\,du.
 \tag{SBD.16}
\]
The second equality moves the finite transvection past \(s\)
as an element of \(K\); finite and real factors commute.
All integrals are ordinary convergent integrals: the integrand
is continuous on a set with compact closure.  The volume of
\(B\) is positive and finite, so division proves (SBD.15).
Changing either Haar normalization multiplies the zero integral
by a nonzero constant.  This auxiliary transvection measure
does not change the ambient target measure chosen in FV11.
% END SBD8

% BEGIN SBD9
\paragraph{SBD9: every proper rational parabolic has a central original transvection.}
\label{claim:SBD9}
Let \(P\) be any proper \(\mathbb Q\)-parabolic of the actual
\(\mathbf G=\operatorname{Res}_{F/\mathbb Q}U(J)\), and let
\(N=R_u(P)\).  There is a nonzero original isotropic
\(\xi\in E^n\) such that the subgroup
\(\operatorname{Res}_{F/\mathbb Q}Z_\xi\) from SBD7 lies in
the center of \(N\).  This assertion covers every rational
parabolic appearing in the classical cusp condition.

\emph{Proof and checked published input.}
We use J.~S.~Milne, \emph{Algebraic Groups} (2017), Theorem~25.1
(printed p.~544) and Proposition~17.60 (printed p.~372), in the
author's \texttt{i\allowbreak AG2017.\allowbreak pdf}.  For a reductive group over its
base field these state that a parabolic is \(P(\lambda)\) for a
cocharacter over that field, and that its unipotent radical is
the subgroup \(U(\lambda)\) contracted to the identity by that
cocharacter.  FV1--FV4 give a smooth connected reductive
\(\mathbb Q\)-group here; characteristic zero and the indicated
rational definition of \(P\) check all the hypotheses.  No
splitness of the unitary group over \(F\) is assumed.

By the adjunction defining restriction of scalars, a
\(\mathbb Q\)-cocharacter of \(\mathbf G\) corresponds to an
\(F\)-cocharacter of \(U(J)\).  On \(E^n\) write its weight
decomposition as \(\bigoplus_m V_m\).  These are actual
\(E\)-subspaces: the cocharacter acts \(E\)-linearly and the
split torus has integral weights.  Hermitian invariance implies
\(h(V_m,V_l)=0\) unless \(m+l=0\), because the scalar
parameter lies in \(F\) and is fixed by \(\sigma\).
Nonsingularity pairs \(V_m\) perfectly with \(V_{-m}\).
The properness of \(P\) rules out all weights being zero.
Consequently the largest weight \(M\) is positive and any
nonzero \(\xi\in V_M\) is isotropic.

The contraction condition for \(u\in U(\lambda)\) says that
\(u-1\) strictly raises weights.  Indeed the block from
\(V_j\) to \(V_i\) is multiplied by \(t^{i-j}\), so a block
with \(i\le j\) must be zero in \(u-1\).  It follows that
\(u\xi=\xi\).  Unitarity then gives
\(\xi^*Ju=\xi^*J\), so \(uT_\xi=T_\xi u=T_\xi\).
The operator \(T_\xi\) maps \(V_{-M}\) to \(V_M\) and
kills the other weight spaces.  Its positive weight \(2M\)
shows that \(1+\eta aT_\xi\) is in \(U(\lambda)\), and
the preceding identities show that it commutes with every
element of that group.  These are polynomial matrix identities
over \(F\), hence apply after restriction of scalars and at
every real or adelic point.  No additional parabolic dictionary
is left as an input.
% END SBD9

% BEGIN SBD10
\paragraph{SBD10: compact unipotent quotients for the actual arithmetic subgroup.}
\label{claim:SBD10}
Let \(N\) be a unipotent \(\mathbb Q\)-subgroup of the original
faithful rational matrix group, and \(\Gamma'\) any arithmetic
subgroup of that group.  Then
\[
 (\Gamma'\cap N(\mathbb R))\backslash N(\mathbb R)
 \quad\hbox{is compact and has a finite nonzero invariant measure.}
 \tag{SBD.17}
\]
This applies to the exact \(\Gamma_s\) in SBD5.

\emph{Proof.}
Use the lattice in the rational representation with whose
stabilizer \(\Gamma'\) is commensurate.  Its intersection with
\(N\) is a discrete integral matrix group, commensurate with
\(\Gamma'\cap N\).  The Lie algebra \(\mathfrak n\) over
\(\mathbb Q\) is nilpotent and has a rational flag
\[
 0=\mathfrak n_0\subset\mathfrak n_1\subset\cdots
  \subset\mathfrak n_r=\mathfrak n,\qquad
 \dim\mathfrak n_j=j,\qquad
 [\mathfrak n,\mathfrak n_j]\subset\mathfrak n_{j-1}.
 \tag{SBD.18}
\]
To construct it, choose a nonzero rational vector in the center
(the kernel of rational linear bracket equations), quotient by
its span, and induct on dimension.  Choose an adapted rational
basis \(X_1,\ldots,X_r\).

For unipotent matrix groups in characteristic zero the finite
matrix polynomials \(\exp\) and \(\log\) are inverse between
the Lie algebra and the group.  Here is a direct check of the
point that can matter in this use.  If \(u\in N(\mathbb R)\),
the polynomial curve \(\exp(t\log u)\) equals \(u^k\) at
every nonnegative integer \(k\); each defining polynomial of
\(N\) therefore vanishes on the whole curve, and its tangent
\(\log u\) lies in \(\mathfrak n(\mathbb R)\).  Conversely
the flow of the left-invariant tangent field of
\(X\in\mathfrak n(\mathbb R)\) stays in the closed Lie
subgroup \(N(\mathbb R)\) and is the matrix exponential.
The finite exponential and logarithm identities now prove the
claim.  In particular the group is connected and simply
connected.  Using (SBD.18) successively gives unique coordinates
\[
 (t_r,\ldots,t_1)\longmapsto
        \exp(t_rX_r)\cdots\exp(t_1X_1).
 \tag{SBD.19}
\]
Uniqueness follows by taking the one-dimensional quotient by
\(\mathfrak n_{r-1}\), then descending the flag; the same
argument proves existence.  These coordinates are polynomial
with polynomial inverse, by the finite multiplication and
logarithm formulas.

For each rational nilpotent matrix \(X_j\), choose an integer
\(d_j>0\) clearing the finitely many denominators in
\(\exp(\pm d_jX_j)\), so both matrices preserve the chosen
integral lattice.  Finite index then puts a positive power of
\(\exp(d_jX_j)\) in \(\Gamma'\cap N\).  A common multiple
\(b>0\) gives \(\exp(bX_j)\in\Gamma'\cap N\) for every
\(j\).  Left multiplication by \(\exp(-kbX_j)\) changes the
\(j\)-th coordinate in (SBD.19) by \(-kb\), leaves every
higher coordinate unchanged, and may change lower coordinates.
This follows from the bracket inclusion in (SBD.18).
Reduce coordinates in the order \(r,r-1,\ldots,1\).
Every left coset therefore meets the compact image of
\([0,b]^r\).  Discreteness makes the quotient Hausdorff, so
it is compact.

The triangular coordinates also make left and right translation
Jacobian one; equivalently \(\operatorname{ad}X\) is nilpotent
and has trace zero for every \(X\).  Thus \(N(\mathbb R)\)
is unimodular.  Haar measure descends to the discrete left
quotient.  A finite cover by local sections of its compact
quotient proves finiteness, and any nonempty local section
proves positivity.  Normalize it to probability when convenient.
% END SBD10

% BEGIN SBD11
\paragraph{SBD11: the exact classical cusp hypothesis is satisfied.}
\label{claim:SBD11}
Under the hypotheses of SBD8, each \(F_s\) is cuspidal for its
actual arithmetic group \(\Gamma_s\): for every proper
\(\mathbb Q\)-parabolic of \(\mathbf G\), with real unipotent
radical \(N(\mathbb R)\),
\[
 \int_{(\Gamma_s\cap N(\mathbb R))\backslash N(\mathbb R)}
          F_s(na)\,dn=0\qquad(a\in G_\infty).
 \tag{SBD.20}
\]
In particular the hypothesis is not being inferred merely
from a name such as \texttt{Is\allowbreak Cusp\allowbreak Fun}.

\emph{Proof.}
Choose the central transvection subgroup \(Z\subset N\)
provided by SBD9.  The quotients by \(\Gamma_s\cap N\) and
\(\Gamma_s\cap Z\) are compact by SBD10 and SBD6 respectively;
give each its invariant probability measure.  For fixed \(n\),
SBD8 at \(na\) gives
\(\int_{(\Gamma_s\cap Z)\backslash Z}F_s(zna)\,dz=0\).
Integrate in \(n\) and interchange the two compact integrals.
For fixed \(z\), centrality makes left translation by \(z\)
on the \(N\)-quotient equal to right translation, which
preserves its invariant measure.  The inner integral is
therefore independent of \(z\) and equals the left side of
(SBD.20).  Its average is zero, proving the assertion.
Continuity on compact quotients justifies Fubini with genuine
integrability.  If the group has no proper rational parabolics,
the cusp requirement is empty, exactly as in the published
theorem.
% END SBD11

% BEGIN SBD12
\paragraph{SBD12: checked rapid decay gives a global bound on each real component.}
\label{claim:SBD12}
Let \(\phi\) be the initial extremal function of SBD2.  Suppose
it is an original automorphic function, has the original
uniform moderate growth, and satisfies original \(\mathrm{ZCusp}\).
For every fixed \(s\in G_f\),
\[
                   \sup_{a\in G_\infty}|\phi(sa)|<\infty.
 \tag{SBD.21}
\]

\emph{Checked published application and proof.}
Use S.~D.~Miller and W.~Schmid, \emph{On the rapid decay of
cuspidal automorphic forms}, Advances in Mathematics 231 (2012),
940--964.  Numbering here is that of the inspected author
preprint \texttt{ar\allowbreak Xiv:\allowbreak 1106.\allowbreak 2149v1}: Theorem~A (printed p.~3),
the finite Siegel cover (2.8) (p.~8), and the precise
automorphic and cusp hypotheses (6.1), (6.3) (p.~29).
Theorem~A bounds a smooth cuspidal automorphic form of moderate
growth by any negative power of a faithful linear norm on each
Siegel set.  The definition includes finiteness under the full
enveloping-algebra center.

Here is the complete application check.  FV1--FV4 identify the
group as real points of a reductive \(\mathbb Q\)-group.  Its
real center is the product of the scalar circles in the original
unitary factors, hence compact.  FV10 and SBD5 give the actual
arithmetic subgroup, smoothness, and moderate growth in the
faithful norm (SBD.10).  SBD3--SBD4 give full center finiteness
on the initial function, before using any target Hilbert space.
SBD11 proves the zero integral for \emph{every} proper rational
parabolic, with convergent classical integrals and the same
arithmetic subgroup.  These are all the published hypotheses;
no \(K\)-finiteness or Hecke-finiteness input is being silently
added.

One still must pass from a Siegel set to the whole group.
The finite cover (2.8) gives a finite set of rational elements
\(c\) and a Siegel set \(\mathfrak S\) with
\(G_\infty=\bigcup_c\Gamma_s c\mathfrak S\).
For each such \(c\), the function \(a\mapsto F_s(ca)\)
has arithmetic group \(c^{-1}\Gamma_s c\).  It is smooth;
its center finiteness is unchanged because right differential
operators commute with left translation.  Its growth follows
from \(\|ca\|_{\mathrm{lin}}\le
\|c\|_{\mathrm{lin}}\|a\|_{\mathrm{lin}}\).
Conjugation carries its rational parabolics and unipotent
quotients to those for \(F_s\), so its cusp integrals are
zero too.  Apply Theorem~A to this translated function on
\(\mathfrak S\), with exponent one.  Since
\(\|a\|_{\mathrm{lin}}\ge1\), it is bounded there.
The maximum of the finitely many constants bounds \(F_s\)
on the whole group.  This also treats the anisotropic case
of the finite Siegel cover; rank zero was handled in SBD5.

All uses of the published result are thus direct applications
after the hypotheses have been checked.  Its proof is a deep
published input, not a newly reconstructed elementary argument.
% END SBD12

% BEGIN SBD13
\paragraph{SBD13: the initial seed bound, with the original quantifier order.}
\label{claim:SBD13}
Fix the same original \(R\), with its original level and model
laws, Weil datum, vacuum, theta-polynomial law and one
simultaneous \(io_i\) before every parameter \(t\).
Choose the unit extremal tuple \(e_i\) of HCU20 before \(t\)
and \(x\).  Let \(U_{i,t,x}\) be the literal integral in
TAE13, on its original source measure and domain.  Then
\[
 \forall i\ \forall t\ \forall x\in(R.HV(i,t)).H\
 \exists C_{i,t,x}\in\mathbb R\ \forall g\in G(W_i):
       |U_{i,t,x}(e_i)(g)|\le C_{i,t,x}.
 \tag{SBD.22}
\]
No uniform constant in \(t\) or \(x\) is asserted or needed.

\emph{Proof.}
TAE13, CTC13 and FIA8 give, for the same initial integral,
smoothness, rational invariance, original principal-level
invariance, a common derivative-growth exponent, compact
equivariance, the lowering equations, and original
\(\mathrm{ZCusp}\).  HCU13 and HCU20 supply the exact model
and extremal hypotheses of SBD2.  The simultaneous \(io_i\)
is the one supplied by JIO19 on the original premises; it is
not selected anew for the bound.

Choose a compact open finite subgroup \(K_i\) inside the
original principal target level, using the exact open-level
calculation in SPR25 and SBD5.  This choice depends on the
original \(R,i\), not on \(t,x\).  FV9 gives a finite set
\(S_i\subset G_f(W_i)\) such that every target element is
\(\gamma s a k\), with \(\gamma\) rational,
\(s\in S_i\), \(a\in G_\infty\) and \(k\in K_i\).
For each fixed \(i,t,x\), SBD12 bounds the function on each
\(sG_\infty\).  Rational left invariance and right
\(K_i\)-invariance then give (SBD.22) by taking a finite
maximum.  A zero seed causes no difficulty.  The original
growth exponent precedes every derivative word; this proof
does not exchange any of those quantifiers.  It introduces
no new family or coefficient threshold.
% END SBD13

% BEGIN SBD14
\paragraph{SBD14: the full original lowest-type extension no longer has analytic gaps.}
\label{claim:SBD14}
Let the original generic data \(D,\theta,V,M,w,J_0,\chi,U\)
satisfy all the hypotheses of
\texttt{Theta\allowbreak Map.\allowbreak exists\_\allowbreak extension} except its sufficient
global predicate \texttt{Published.\allowbreak Cusp\allowbreak Rapid\allowbreak Decay}.
In particular keep original \(\mathrm{IsInfFam}\),
\(\mathrm{HolShape}\), \(\mathrm{WtIs}\),
\(J_0\subset G_f\), and, for every lowest tuple, original
automorphy, level covariance, smoothness, uniform moderate
growth and \(\mathrm{ZCusp}\), together with compact
equivariance and lowering.  Then there exists an original
\(T:\mathrm{MMap}(M)\) with
\[
       \mathrm{IsMultMap}(\theta,M,J_0,\chi,T),\qquad
       \mathrm{RestrictsTo}(M,T,U).
 \tag{SBD.23}
\]

\emph{Proof.}
HCU13 supplies \(\mathrm{LowSm}\) and \(\mathrm{LowIrr}\),
and choose one original nonzero extremal tuple.  SBD2--SBD12
apply to its initial value.  Its own \texttt{Is\allowbreak Aut\allowbreak Fun} supplies
the compact open level required in SBD5.  FV9 reduces its
adelic domain to finitely many real components, exactly as
in SBD13, and hence proves its global boundedness.
HCU19 supplies its literal \(\mathrm{HolAd}\).  FV8 supplies
an original finite pair for this same target.  These data are
all the inputs to SPR21, including the separate original
\(\mathrm{ZCusp}\) premise.  SPR21 gives an \(\mathrm{IsMultMap}\)
whose value at this extremal tuple is the actual initial
function.  SPR22, using the same compact equivariance and
the original nonzero tuple, extends that equality to every
lowest tuple.  This proves (SBD.23).

The full pointwise smoothness, every ordered derivative,
Sobolev continuity, central finiteness, cusp property, and
finite and real equivariance are the fields already proved
in SPR14--SPR21.  The finite pair chosen here is the same
one throughout that construction.  In particular no
post-Hilbert growth or scalar fact was used to obtain the
initial bound.
% END SBD14

% BEGIN SBD15
\paragraph{SBD15: the original all-parameter theta map is now constructed.}
\label{claim:SBD15}
On the same original \(R\), retain the premises of SPR25:
the original family, level and model laws, Weil data,
vacuum and theta-polynomial identities, and the actual
simultaneous \(io_i\) of JIO19.  Then there is one family
\(\mathrm{thF}\) such that the literal update
\[
                  R_1=\{R\ \mathrm{with}\ \mathrm{thF}:=\mathrm{thF}\}
\]
satisfies both original predicates
\(R_1.\mathrm{thP\_thF\_mem}\) and
\(R_1.\mathrm{thP\_thF}\), for every original \(i,t\).
All original source measures, kernels, frames, models,
levels and the same \(io_i\) are retained.

\emph{Proof.}
FV11 chooses an original \((\mu_i,\mathcal F_i)\) with
\(\mathrm{FinPair}(W_i,\mu_i,\mathcal F_i)\), before all
\(t,x\).  SBD13 supplies exactly SPR25's other remaining
analytic premise.  Insert these two outputs in SPR25.
For clarity, the pointwise extension first agrees with the
actual convergent integral at the original extremal tuple;
SPR22 gives equality on all lowest tensors.  SPR24's
algebraic linear lifting assembles dependence on the entire
source vector space, without assuming injectivity of its
extension map.  The integral is already bilinear, so the
lift is consistent on the actual range.  These are exactly
the original \texttt{Fires\allowbreak At} and multiplicity-space
membership assertions.  Choice over the original indices
\((i,t)\) then gives the indicated family.

The stored weak \(\mathrm{OrZero}\) selections are activated
by the same existence argument of SPR25; this proves their
predicates and does not identify a stored witness with an
independently chosen one.  No geometric reading, rational
test, row identity or new selected \(R\) is created by this
construction.  The universal theorem
\texttt{Published.\allowbreak Cusp\allowbreak Rapid\allowbreak Decay} remains outside the
claims made here; its stronger conclusion is unnecessary
for this exact theta-map endpoint.
% END SBD15

% BEGIN SBD16
\paragraph{SBD16: closure of the selected cuspidal-extension obligation.}
\label{claim:SBD16}
The selected obligation \texttt{BASE\_\allowbreak THETA\_\allowbreak CUSPIDAL\_\allowbreak EXTENSION}
is now supplied on its original common-data premises.  On
the same selected Bundle and same reading as
\texttt{BASE\_\allowbreak THETA\_\allowbreak READINGS}, each actual convergent
lowest-type integral extends to the original full
\(\mathrm{Wall}_i\), with all original smoothness, finite
level, real equivariance, cusp and Sobolev-continuity
requirements.  One common family threshold before the
indices and finite source basis is preserved.

\emph{Proof.}
The selected contract is the restriction of SBD15 to the
original base parameter and to the finitely many source
vectors used in BTH12.  SBD15 is stronger in those indices:
it holds for every \(t\) and every source vector on every
original datum satisfying the already stated common-data
premises.  Thus the parent threshold at which those common
premises hold also suffices for all of these extensions;
there is no new threshold to place after a basis vector.
The restriction to the actual \(\mathrm{Wall}_i\) is the
literal one in BTH4--BTH7 and SPR25, using their original
models and quotient.  Apply it to the maps of SBD15.  The
fixed original source basis is handled simultaneously by
SPR24, not by unrelated existential extensions.

This closes the analytic leaf as a universally supplied
producer on that same datum.  Its neighboring obligations
for common geometric readings, simultaneous rational
tests, actual base rows, and common selected realizations
remain separate inputs.  In particular this result alone
does not give the base theta square, its rational adjoint,
or the final RH/GRH theorem.  The classification of a
published theorem used inside this proof is independent
of whether those other producers have been completed.
% END SBD16

\paragraph{Publication and source record.}
The checked external inputs are recorded as
\texttt{SBD\_\allowbreak PARABOLIC\_\allowbreak COCHARACTER} and
\texttt{SBD\_\allowbreak RAPID\_\allowbreak DECAY} in this component's provenance and
closure files.  The exact original source anchors include:
\begin{itemize}
\item \texttt{Tied\allowbreak Smooth.\allowbreak Rder\allowbreak C/\allowbreak gelf/\allowbreak Is\allowbreak Aut\allowbreak Fun\allowbreak Sm}
and \texttt{Is\allowbreak Smooth\allowbreak Vector.\allowbreak Killed\allowbreak Fn/\allowbreak ZI/\allowbreak ZCusp/\allowbreak Fin\allowbreak Pair};
\item \texttt{Holomorphic\allowbreak Type\allowbreak At\allowbreak Extremal\allowbreak Of\allowbreak KEquivariant\allowbreak Lowering\allowbreak K}
and \texttt{Exists\allowbreak Sigma\allowbreak List};
\item \texttt{Exists\allowbreak Theta\allowbreak Lift\allowbreak Linear\allowbreak Map\allowbreak Cusp\allowbreak Form\allowbreak Multiplicity}.
\end{itemize}
The alternate full-center argument in SBD3 is independent;
it does not claim to reconstruct the separate source theorem
that cyclic Gelfand elements generate the whole center.
\end{document}
